\section{РЕАЛИЗАЦИЯ И ТЕСТИРОВАНИЕ}

\subsection{Общая структура реализации}

Backend-платформа CoachLink реализована как набор Go-сервисов в монорепозитории. Выбор монорепозитория обусловлен упрощением управления общими зависимостями, единой системой сборки и упрощённой координацией изменений между сервисами на этапе активной разработки.

Каждый сервис реализован как самостоятельный Go-модуль со стандартной внутренней структурой: точка входа, конфигурация, HTTP-обработчики, сервисный слой с бизнес-логикой, модели данных и слой доступа к базе данных или внешним клиентам. Общий код, используемый несколькими сервисами, вынесен в разделяемый пакет --- в частности, типы и названия событий NATS JetStream, что обеспечивает согласованность событийных контрактов между публикующими и подписанными сервисами.

Для локального развёртывания используется Docker Compose: все семь сервисов, PostgreSQL, NATS JetStream и Ollama запускаются как контейнеры. Для каждого сервиса предусмотрен отдельный Dockerfile, обеспечивающий воспроизводимую сборку. Stateful-сервисы применяют SQL-миграции при запуске с помощью библиотеки goose, что гарантирует актуальность схемы базы данных без ручного вмешательства. Конфигурация каждого сервиса задаётся через переменные окружения, что упрощает развёртывание в различных средах. Краткое описание назначения каждого сервиса и его взаимодействий приведено в приложении А.

\subsection{API Gateway}

API Gateway реализован как отдельный Go-сервис на основе фреймворка Echo. Проксирование запросов выполняется через \passthrough{\lstinline!net/http/httputil.ReverseProxy!} из стандартной библиотеки Go. Для каждого backend-сервиса создаётся отдельный экземпляр reverse proxy с пользовательским \passthrough{\lstinline!Director!}, который сохраняет полный путь запроса и подменяет только scheme и host целевого сервиса. При ошибке проксирования пользовательский \passthrough{\lstinline!ErrorHandler!} возвращает клиенту ответ 502 в формате JSON.

Входящие запросы проходят через цепочку middleware в фиксированном порядке: восстановление после паники, генерация идентификатора запроса, обработка CORS, структурированное логирование и проверка JWT. Маршрутизация выполняется по префиксам, как показано в таблице~\ref{tab:gateway-routing}.

\begin{longtable}{|>{\raggedright\arraybackslash}p{0.42\columnwidth}|>{\raggedright\arraybackslash}p{0.46\columnwidth}|}
\caption{Маршрутизация API Gateway по префиксам}\label{tab:gateway-routing}\\
\hline
Префикс & Целевой сервис \\
\hline
\endhead
\hline
\endlastfoot
\passthrough{\lstinline!/api/v1/auth/*!} & Auth Service \\ \hline
\passthrough{\lstinline!/api/v1/users/*!} & User Service \\ \hline
\passthrough{\lstinline!/api/v1/connections/*!} & User Service \\ \hline
\passthrough{\lstinline!/api/v1/groups/*!} & User Service \\ \hline
\passthrough{\lstinline!/api/v1/training/*!} & Training Service \\ \hline
\passthrough{\lstinline!/api/v1/notifications/*!} & Notification Service \\ \hline
\passthrough{\lstinline!/api/v1/analytics/*!} & Analytics Service \\ \hline
\passthrough{\lstinline!/api/v1/ai/*!} & AI Service \\ \hline
\end{longtable}

Gateway выполняет JWT-проверку для всех защищённых маршрутов. Исключениями являются регистрация, вход, обновление токена и health-check. После успешной проверки токена Gateway извлекает из полей токена идентификатор пользователя, роль и логин. Эти значения передаются дальше через заголовки, представленные на рисунке~\ref{src:gateway-headers}.

\begin{figure}
\begin{lstlisting}
X-User-ID: <user_id>
X-User-Role: <coach|athlete>
X-User-Login: <login>
\end{lstlisting}
\caption{Заголовки запроса, передаваемые Gateway внутренним сервисам}
\label{src:gateway-headers}
\end{figure}

Такой механизм позволяет внутренним сервисам принимать решения о доступе без повторной проверки JWT. Например, Training Service проверяет, может ли текущий пользователь создать план, отправить отчёт или получить назначение.

Gateway также отвечает за обработку CORS-заголовков, необходимых для корректной работы web-клиентов, и предоставляет health-check endpoint для проверки доступности сервиса.

\subsection{Auth Service}

Auth Service отвечает за регистрацию, вход и управление токенами. Сервис использует PostgreSQL-базу \passthrough{\lstinline!auth\_db!} и таблицы \passthrough{\lstinline!users!} и \passthrough{\lstinline!refresh\_tokens!}.

\subsubsection{Регистрация}

При регистрации пользователь передаёт логин, email, пароль, ФИО и роль. DTO запроса содержит ограничения:

\begin{itemize}
\tightlist
\item
  логин обязателен, длина от 3 до 50 символов;
\item
  email должен иметь корректный формат;
\item
  пароль обязателен, длина от 8 до 128 символов;
\item
  ФИО обязательно;
\item
  роль должна быть \passthrough{\lstinline!coach!} или \passthrough{\lstinline!athlete!}.
\end{itemize}

Дополнительно сервис проверяет формат логина регулярным выражением (рисунок~\ref{src:login-regex}).

\begin{figure}
\begin{lstlisting}
^[a-zA-Z0-9-]+$
\end{lstlisting}
\caption{Регулярное выражение для проверки формата логина}
\label{src:login-regex}
\end{figure}

Это ограничение исключает пробелы, кириллицу и специальные символы в логине, что упрощает поиск и использование логина в API.

Пароль хешируется через bcrypt с cost-фактором 12, что обеспечивает баланс между безопасностью и временем вычисления хеша. В базе хранится только \passthrough{\lstinline!password\_hash!}, исходный пароль не сохраняется. Значение cost-фактора задаётся через переменную окружения и может быть изменено без модификации кода. После создания пользователя сервис публикует событие \passthrough{\lstinline!coachlink.user.registered!}, содержащее идентификатор пользователя, логин, ФИО, email и роль.

\subsubsection{Вход и токены}

При входе сервис ищет пользователя по логину и сравнивает пароль с bcrypt-хешем. При успешной проверке создаётся пара токенов:

\begin{itemize}
\tightlist
\item
  access token --- JWT, подписанный алгоритмом HMAC-SHA256, со сроком жизни 15 минут;
\item
  refresh token --- случайный UUID со сроком жизни 30 дней.
\end{itemize}

Структура полей access token описана в разделе проектирования (рисунок~\ref{src:design-jwt-payload}).

Refresh token не хранится в открытом виде. Сервис вычисляет SHA-256 хеш токена и сохраняет только хеш. Такой подход снижает риск компрометации при утечке базы данных: даже получив содержимое таблицы, злоумышленник не сможет восстановить исходный токен. При refresh старый хеш удаляется из базы, затем создаётся новая пара access/refresh. При logout refresh token также удаляется, что делает невозможным дальнейшее обновление сессии. Auth Service предоставляет четыре публичных endpoint: регистрацию, вход, обновление токена и выход.

\subsection{User Service}

User Service хранит пользовательские профили, связи тренер-спортсмен и тренировочные группы. Он использует базу \passthrough{\lstinline!user\_db!}.

\subsubsection{Синхронизация профилей}

Профиль пользователя создаётся не напрямую при регистрации, а через событие \passthrough{\lstinline!coachlink.user.registered!}. Auth Service публикует событие после успешной регистрации, а User Service подписан на него через NATS JetStream. При получении события User Service создаёт запись в таблице \passthrough{\lstinline!user\_profiles!}.

Такое разделение позволяет Auth Service отвечать только за учётные данные и токены, а User Service --- за доменную информацию, используемую в тренировочном процессе.

\subsubsection{Поиск пользователей}

Сервис предоставляет поиск по ФИО и логину. Поиск используется, например, когда спортсмен ищет тренера для отправки заявки или тренер ищет спортсмена. Минимальная длина поискового запроса --- два символа. Это ограничение снижает нагрузку на индекс при поиске по коротким строкам.

Для ускорения поиска в базе включено расширение \passthrough{\lstinline!pg\_trgm!}, а для полей имени и логина создаются trigram-индексы. Это подходит для поиска подстрок и автодополнения.

\subsubsection{Заявки и связи}

Связь тренера и спортсмена создаётся через заявку. Основной сценарий:

\begin{enumerate}
\def\labelenumi{\arabic{enumi}.}
\tightlist
\item
  спортсмен находит тренера;
\item
  спортсмен отправляет заявку;
\item
  тренер получает уведомление;
\item
  тренер принимает или отклоняет заявку;
\item
  при принятии создаётся запись в \passthrough{\lstinline!coach\_athlete\_relations!}.
\end{enumerate}

В первой версии у спортсмена может быть один тренер. Это ограничение задаётся уникальностью \passthrough{\lstinline!athlete\_id!} в таблице связей. Такое решение упрощает логику доступа: Training Service и Analytics Service могут однозначно определять, какой тренер отвечает за спортсмена.

\subsubsection{Тренировочные группы}

Тренер может создавать тренировочные группы и добавлять туда своих спортсменов. Группа используется для массового назначения тренировочного плана: при создании плана с указанием \passthrough{\lstinline!group\_id!} Training Service запрашивает состав группы через внутренний API User Service и формирует назначения для каждого участника.

Перед добавлением спортсмена в группу сервис проверяет наличие записи в таблице \passthrough{\lstinline!coach\_athlete\_relations!}. Если связь отсутствует, запрос отклоняется. При удалении группы записи в таблице участников удаляются каскадно. При добавлении или удалении спортсмена из группы публикуются события \passthrough{\lstinline!coachlink.group.athlete\_added!} и \passthrough{\lstinline!coachlink.group.athlete\_removed!}, на которые реагирует Notification Service.

\subsubsection{Внутренние API}

User Service предоставляет внутренние API для других сервисов (таблица~\ref{tab:user-internal-api}).

\begin{longtable}[]{|
  >{\raggedright\arraybackslash}p{(\columnwidth - 2\tabcolsep) * \real{0.4700}}|
  >{\raggedright\arraybackslash}p{(\columnwidth - 2\tabcolsep) * \real{0.4700}}|}
\caption{Внутренние API User Service}\label{tab:user-internal-api}\\
\hline
\begin{minipage}[b]{\linewidth}\raggedright
Endpoint
\end{minipage} & \begin{minipage}[b]{\linewidth}\raggedright
Использование
\end{minipage} \\
\hline
\endfirsthead
\caption*{Продолжение таблицы \ref{tab:user-internal-api}}\\
\hline
\begin{minipage}[b]{\linewidth}\raggedright
Endpoint
\end{minipage} & \begin{minipage}[b]{\linewidth}\raggedright
Использование
\end{minipage} \\
\hline
\endhead
\hline
\endlastfoot
GET \path|/internal/users/{userId}| & Получение профиля пользователя \\ \hline
GET \path|/internal/groups/{groupId}/members| & Получение состава группы \\ \hline
GET \path|/internal/coach/{coachId}/has-athlete/{athleteId}| & Проверка связи тренера и спортсмена \\ \hline
\end{longtable}

Эти API не предназначены для внешних клиентов. Они позволяют Training Service, Analytics Service и AI Service выполнять предметные проверки и получать необходимые данные.

\subsection{Training Service}

Training Service является центральным сервисом тренировочного процесса. Он использует базу \passthrough{\lstinline!training\_db!}.

\subsubsection{Тренировочные планы}

Тренировочный план содержит название, описание, дату и идентификатор тренера. План может быть назначен:

\begin{itemize}
\tightlist
\item
  одному или нескольким спортсменам по списку \passthrough{\lstinline!athlete\_ids!};
\item
  всем спортсменам группы через \passthrough{\lstinline!group\_id!};
\item
  комбинации группы и отдельных спортсменов.
\end{itemize}

Если указан \passthrough{\lstinline!group\_id!}, Training Service обращается к User Service и получает участников группы. Участники группы и спортсмены из \passthrough{\lstinline!athlete\_ids!} объединяются через ассоциативный массив с дедупликацией по идентификатору: если спортсмен указан одновременно в группе и в списке, он получит только одно назначение. Затем создаётся один тренировочный план и набор назначений для каждого спортсмена. При этом используется стратегия частичного успеха: если создание назначения для одного из спортсменов завершается ошибкой, остальные назначения продолжают создаваться.

\subsubsection{Назначения}

Назначение представляет связь между тренировочным планом и конкретным спортсменом. В назначении хранятся:

\begin{itemize}
\tightlist
\item
  \passthrough{\lstinline!plan\_id!};
\item
  \passthrough{\lstinline!athlete\_id!};
\item
  \passthrough{\lstinline!coach\_id!};
\item
  имя и логин спортсмена;
\item
  имя и логин тренера;
\item
  статус;
\item
  время назначения;
\item
  время выполнения;
\item
  время архивации.
\end{itemize}

Денормализация имени и логина спортсмена используется для фильтрации списков назначений без дополнительного запроса в User Service.

Статусы назначений приведены в таблице~\ref{tab:assignment-statuses}.

\begin{longtable}{|>{\raggedright\arraybackslash}p{0.42\columnwidth}|>{\raggedright\arraybackslash}p{0.46\columnwidth}|}
\caption{Статусы тренировочных назначений}\label{tab:assignment-statuses}\\
\hline
Статус & Значение \\
\hline
\endhead
\hline
\endlastfoot
\passthrough{\lstinline!assigned!} & Задание назначено и ожидает отчёта \\ \hline
\passthrough{\lstinline!completed!} & Спортсмен отправил отчёт \\ \hline
\passthrough{\lstinline!archived!} & Тренер перенёс выполненное назначение в архив \\ \hline
\end{longtable}

\subsubsection{Просроченность задания}

Просроченность вычисляется функцией \passthrough{\lstinline!ComputeIsOverdue!}. Задание считается просроченным, если:

\begin{itemize}
\tightlist
\item
  статус равен \passthrough{\lstinline!assigned!};
\item
  дата тренировки уже прошла.
\end{itemize}

Для определения просроченности используется grace-период: задание считается просроченным не в момент наступления даты тренировки, а спустя 24 часа после неё. Это даёт спортсмену возможность отправить отчёт в течение дня после тренировки. Просроченность не хранится в базе данных как отдельное поле, а вычисляется при каждом чтении. Это исключает необходимость фоновых задач и ситуацию, когда сохранённое значение устаревает.

\subsubsection{Отчёты}

Перед созданием отчёта выполняется двойная проверка: сервис убеждается, что назначение находится в статусе \passthrough{\lstinline!assigned!}, и что для данного назначения ещё не существует отчёта. Если отчёт уже был отправлен ранее, возвращается ошибка 409 Conflict. Такая двойная защита предотвращает дублирование отчётов даже при повторных запросах. После успешного создания отчёта статус назначения меняется на \passthrough{\lstinline!completed!}, а в поле \passthrough{\lstinline!completed\_at!} записывается текущее время.

Формат отчёта представлен на рисунке~\ref{src:report-format}.

\begin{figure}
\begin{lstlisting}
{
  "content": "Тренировка прошла хорошо, последние интервалы дались тяжело",
  "duration_minutes": 75,
  "perceived_effort": 7,
  "max_heart_rate": 188,
  "avg_heart_rate": 162,
  "distance_km": 12.4
}
\end{lstlisting}
\caption{Структура JSON-тела запроса на отправку отчёта}
\label{src:report-format}
\end{figure}

Поля отчёта выбраны с учётом первой реализации для лёгкой атлетики. Они позволяют фиксировать как объективные показатели, так и субъективную оценку нагрузки.

\subsubsection{Шаблоны тренировок}

Шаблоны тренировок позволяют тренеру сохранять часто используемые описания тренировок для повторного применения. Шаблон хранит те же предметные поля, что и план --- название и описание, --- но не привязан к дате и спортсменам. При создании тренировочного плана тренер может указать флаг \passthrough{\lstinline!save\_as\_template!}, и система автоматически сохранит копию описания как шаблон. В дальнейшем при создании нового плана тренер может выбрать существующий шаблон, и его поля будут использованы как основа. Каждый шаблон принадлежит конкретному тренеру и недоступен другим пользователям, что обеспечивается проверкой \passthrough{\lstinline!coach\_id!} на уровне сервисного слоя.

\subsubsection{События Training Service}

Training Service публикует события, перечисленные в таблице~\ref{tab:training-service-events}.

\begin{longtable}{|>{\raggedright\arraybackslash}p{0.42\columnwidth}|>{\raggedright\arraybackslash}p{0.46\columnwidth}|}
\caption{События Training Service}\label{tab:training-service-events}\\
\hline
Событие & Когда публикуется \\
\hline
\endhead
\hline
\endlastfoot
\passthrough{\lstinline!coachlink.training.assigned!} & Создано назначение \\ \hline
\passthrough{\lstinline!coachlink.training.deleted!} & Назначение удалено \\ \hline
\passthrough{\lstinline!coachlink.report.submitted!} & Спортсмен отправил отчёт \\ \hline
\end{longtable}

Эти события используются Notification Service для создания уведомлений.

\subsection{Notification Service}

Notification Service отвечает за создание, хранение и предоставление уведомлений пользователям. Сервис подписан на события NATS JetStream и реагирует на них созданием записей в базе \passthrough{\lstinline!notification\_db!}. Благодаря событийной модели сервисы-источники событий не зависят от способа доставки уведомлений: Training Service публикует событие о назначении тренировки, а Notification Service самостоятельно решает, как уведомить спортсмена.

\subsubsection{Обработка событий}

Для каждого типа события создаётся отдельный durable consumer в NATS JetStream. Durable consumer сохраняет позицию чтения, что позволяет сервису продолжить обработку после перезапуска и не потерять события, опубликованные во время простоя.

Примеры соответствия событий и уведомлений представлены в таблице~\ref{tab:event-notification-map}.

\begin{longtable}{|>{\raggedright\arraybackslash}p{0.38\columnwidth}|>{\raggedright\arraybackslash}p{0.18\columnwidth}|>{\raggedright\arraybackslash}p{0.32\columnwidth}|}
\caption{Соответствие событий и создаваемых уведомлений}\label{tab:event-notification-map}\\
\hline
Событие & Получатель & Тип уведомления \\
\hline
\endhead
\hline
\endlastfoot
\passthrough{\lstinline!connection.requested!} & Тренер & Новая заявка от спортсмена \\ \hline
\passthrough{\lstinline!connection.accepted!} & Спортсмен & Заявка принята \\ \hline
\passthrough{\lstinline!training.assigned!} & Спортсмен & Новое тренировочное задание \\ \hline
\passthrough{\lstinline!report.submitted!} & Тренер & Новый отчёт \\ \hline
\passthrough{\lstinline!group.athlete\_added!} & Спортсмен & Добавление в группу \\ \hline
\end{longtable}

Если уведомление успешно создано, сообщение подтверждается через \passthrough{\lstinline!Ack!}. Если произошла ошибка, сообщение отклоняется через \passthrough{\lstinline!Nak!}, чтобы оно могло быть обработано повторно.

\subsubsection{API уведомлений}

Notification Service предоставляет:

\begin{itemize}
\tightlist
\item
  получение списка уведомлений с пагинацией;
\item
  фильтрацию по признаку прочтения;
\item
  получение количества непрочитанных уведомлений;
\item
  пометку одного уведомления прочитанным;
\item
  пометку всех уведомлений прочитанными;
\item
  регистрацию FCM-токена устройства.
\end{itemize}

FCM-интеграция является опциональной: если credentials не настроены, push-доставка отключается, но серверные уведомления продолжают сохраняться в базе данных и доступны через API. Это позволяет вести локальную разработку и тестирование без внешних Firebase-зависимостей. При наличии FCM-конфигурации push-уведомления отправляются в отдельной горутине и не блокируют создание серверного уведомления. В push-данные включаются тип события и идентификатор уведомления, что позволяет мобильному клиенту выполнять маршрутизацию к нужному экрану при нажатии на push. Для отправки на несколько устройств одного пользователя используется метод \passthrough{\lstinline!SendEachForMulticast!} из Firebase Admin SDK.

\subsection{Analytics Service}

Analytics Service реализован как stateless-сервис без собственной базы данных. При каждом запросе он обращается к Training Service через внутренние HTTP API, получает необходимые данные и выполняет агрегацию при обработке запроса. Такой подход исключает проблему рассинхронизации аналитической базы с основной базой тренировок: результат всегда вычисляется по актуальным данным. Сервис предоставляет три основных представления: сводку по спортсмену, прогресс по периодам и обзор для тренера.

\subsubsection{Сводка спортсмена}

Сводка представляет собой агрегированный набор показателей за весь период тренировок спортсмена. Она позволяет тренеру оценить общий объём работы, регулярность выполнения заданий и средние значения ключевых метрик. Сводка включает:

\begin{itemize}
\tightlist
\item
  количество отчётов;
\item
  суммарную длительность;
\item
  среднюю длительность;
\item
  суммарную дистанцию;
\item
  среднюю воспринимаемую нагрузку;
\item
  средний пульс;
\item
  максимальный пульс;
\item
  количество назначений;
\item
  процент выполнения.
\end{itemize}

\subsubsection{Прогресс по периодам}

Помимо общей сводки, тренеру важно отслеживать изменения показателей во времени. Analytics Service группирует отчёты по неделям или месяцам в зависимости от параметра запроса. Для каждого периода рассчитываются:

\begin{itemize}
\tightlist
\item
  количество отчётов;
\item
  суммарная длительность;
\item
  средняя нагрузка;
\item
  средний пульс;
\item
  суммарная дистанция.
\end{itemize}

Такая группировка позволяет выявлять тенденции в нагрузке и сравнивать периоды между собой, а не ограничиваться итоговыми значениями за всё время.

\subsubsection{Обзор тренера}

Для тренера сервис формирует обзор по всем его спортсменам: общее количество спортсменов, суммарное число назначений и отчётов, а также средний процент выполнения заданий. Этот endpoint предназначен для экрана сводки в клиентском приложении, где тренер может быстро оценить общую картину по своей группе, не просматривая статистику каждого спортсмена отдельно.

\subsection{AI Service}

AI Service предоставляет публичный endpoint \lstinline!POST /api/v1/ai/athletes/{id}/recommendations! для генерации рекомендаций.

\subsubsection{Проверка доступа}

Сервис доступен только пользователю с ролью \passthrough{\lstinline!coach!}. После получения запроса AI Service обращается к User Service и проверяет, связан ли указанный спортсмен с текущим тренером. Если связи нет, возвращается ошибка доступа.

\subsubsection{Получение данных}

После проверки доступа сервис получает из Analytics Service:

\begin{itemize}
\tightlist
\item
  сводную статистику спортсмена;
\item
  последние отчёты спортсмена.
\end{itemize}

Если отчётов нет, сервис возвращает ошибку недостатка данных. Это необходимо, потому что AI-рекомендация должна опираться на фактическую тренировочную историю, а не на пустой контекст.

\subsubsection{Формирование prompt}

Prompt состоит из двух частей: системного и пользовательского. Системный prompt задаёт роль LLM (помощник тренера по лёгкой атлетике) и требует строгий формат ответа из трёх секций: <<Тенденции>> --- наблюдения с конкретными цифрами из данных, <<Риски>> --- сигналы тревоги, <<Рекомендации>> --- конкретные действия на ближайшие 1--2 недели. Правила запрещают вводные фразы, вопросы к тренеру и утверждения без числового подкрепления из отчётов.

Пользовательский prompt формируется динамически и содержит общую статистику спортсмена, последние тренировки с длительностью, дистанцией, субъективной нагрузкой, показателями пульса и комментариями спортсмена, а также дополнительный контекст от тренера, если он был передан в запросе.

Количество отчётов ограничивается пятью. Это снижает размер prompt и ускоряет генерацию, сохраняя при этом достаточный контекст для анализа тенденций.

\subsubsection{Генерация через Ollama}

AI Service обращается к Ollama API (рисунок~\ref{src:ollama-endpoint}).

\begin{figure}
\begin{lstlisting}
POST /api/chat
\end{lstlisting}
\caption{Endpoint Ollama API для генерации текста}
\label{src:ollama-endpoint}
\end{figure}

Используется модель \passthrough{\lstinline!gemma3:4b!} (Google Gemma 3, 4 млрд параметров). Для управления генерацией заданы параметры: \passthrough{\lstinline!temperature: 0.7!} (баланс между детерминированностью и вариативностью), \passthrough{\lstinline!num\_ctx: 2048!} (размер контекстного окна в токенах) и \passthrough{\lstinline!num\_predict: 1200!} (максимальная длина ответа). HTTP-таймаут вызова к Ollama составляет 5 минут, что учитывает время генерации на CPU. Ответ возвращается тренеру в виде структуры, представленной на рисунке~\ref{src:ai-response}.

\begin{figure}
\begin{lstlisting}
{
  "athlete_id": "uuid",
  "type": "recommendations",
  "content": "Текст рекомендаций от LLM...",
  "generated_at": "2026-04-13T12:00:00Z",
  "model": "gemma3:4b"
}
\end{lstlisting}
\caption{Структура JSON-ответа AI Service}
\label{src:ai-response}
\end{figure}

AI-модуль не заменяет тренера, а предоставляет вспомогательный текст, который тренер может учитывать при принятии решения о корректировке нагрузки. Примеры тел HTTP-запросов и ответов для основных сценариев (отправка отчёта, получение AI-рекомендации) приведены в приложении В.

\subsection{Тестирование}

\subsubsection{Уровни тестирования}

Для проверки корректности backend-платформы применяется многоуровневый подход к тестированию, который позволяет проверять систему на разных уровнях абстракции.

Unit-тесты проверяют бизнес-логику отдельных сервисов в изоляции от инфраструктуры. Зависимости заменяются mock-объектами, что позволяет быстро проверять правила доступа, обработку ошибок, граничные условия и вычисления без необходимости запуска баз данных или брокера сообщений.

E2E smoke-тест проверяет основной пользовательский сценарий через HTTP API при полностью поднятой платформе. Он проходит путь от регистрации до получения AI-рекомендации и служит для быстрой проверки того, что все сервисы запускаются и взаимодействуют корректно.

Интеграционные тесты проверяют группы endpoint через API Gateway. Они наиболее приближены к поведению реального внешнего клиента, потому что запросы проходят через тот же входной слой, включая JWT-проверку и маршрутизацию, что и мобильное приложение. Каждый запуск использует уникальные логины, что позволяет выполнять тесты повторно без очистки базы данных.

\subsubsection{Unit-тесты}

Unit-тесты написаны для шести сервисов и используют табличный стиль (table-driven tests) с библиотекой testify; зависимости (репозитории, NATS-публикаторы, HTTP-клиенты) заменяются mock-объектами. Тесты сервисного слоя проверяют бизнес-правила в изоляции, а тесты слоя обработчиков (auth-service, user-service) дополнительно прогоняют полный HTTP-путь через Echo с помощью \passthrough{\lstinline!net/http/httptest!}: разбор запроса, валидацию, вызов сервиса и формирование кодов ответа. Всего реализовано 112 unit-тестов; их распределение по сервисам приведено в таблице~\ref{tab:unit-test-coverage}.

\begin{longtable}[]{|
  >{\raggedright\arraybackslash}p{(\columnwidth - 4\tabcolsep) * \real{0.6000}}|
  >{\centering\arraybackslash}p{(\columnwidth - 4\tabcolsep) * \real{0.3000}}|}
\caption{Распределение unit-тестов по сервисам}\label{tab:unit-test-coverage}\\
\hline
Сервис & Число тестов \\
\hline
\endfirsthead
\caption*{Продолжение таблицы \ref{tab:unit-test-coverage}}\\
\hline
Сервис & Число тестов \\
\hline
\endhead
\hline
\endlastfoot
\passthrough{\lstinline!training-service!} & 31 \\ \hline
\passthrough{\lstinline!user-service!} & 28 \\ \hline
\passthrough{\lstinline!auth-service!} & 23 \\ \hline
\passthrough{\lstinline!analytics-service!} & 11 \\ \hline
\passthrough{\lstinline!ai-service!} & 10 \\ \hline
\passthrough{\lstinline!notification-service!} & 9 \\ \hline
\end{longtable}

Содержательно unit-тесты покрывают: для \passthrough{\lstinline!auth-service!} --- валидацию формата логина, регистрацию и дубликаты (409), вход с неверными учётными данными (401), одноразовую ротацию refresh-токена; для \passthrough{\lstinline!user-service!} --- обязательность заголовков аутентификации (401), валидацию поиска и имени группы (400), ролевые ограничения (403); для \passthrough{\lstinline!training-service!} --- вычисление просроченности, контроль доступа к заданиям (403/404), создание плана с назначением каждому из нескольких спортсменов, разворачивание группы в участников и сохранение как шаблон; для \passthrough{\lstinline!notification-service!} --- регистрацию FCM-токена, пагинацию и создание уведомления из события; для \passthrough{\lstinline!analytics-service!} --- агрегацию сводки и группировку по периодам; для \passthrough{\lstinline!ai-service!} --- успешную генерацию, недостаток данных, ошибки Ollama и лимит отчётов в prompt.

Unit-тесты запускаются командой \passthrough{\lstinline!make test-unit!} без Docker и внешних зависимостей. Все 112 тестов шести сервисов выполняются успешно. Тестирование сознательно сосредоточено на слое бизнес-логики (сервисном) и обработчиках: слои доступа к базе данных и конфигурации проверяются на интеграционном уровне, где задействована реальная СУБД.

\subsubsection{E2E и интеграционные тесты}

E2E и интеграционные тесты выполняются при поднятой платформе. E2E smoke-сценарий запускается командой \passthrough{\lstinline!make test-e2e!} и проходит полный пользовательский путь из 19 шагов: от регистрации тренера и спортсмена до отправки отчёта и получения уведомлений. Интеграционные тесты запускаются командой \passthrough{\lstinline!make test-integration!} в отдельном контейнере и проверяют все группы endpoint через API Gateway --- так же, как это делает мобильное приложение, включая JWT-проверку и маршрутизацию. Каждый запуск использует уникальные логины, поэтому тесты можно выполнять повторно без очистки базы данных. Все проверки выполняются успешно.

\subsubsection{Таблица результатов}

Результаты всех тестовых проверок приведены в таблице~\ref{tab:test-results}. Таблица тестирования с критериями фиксации результатов приведена в приложении Е.

\begin{longtable}[]{|
  >{\raggedright\arraybackslash}p{(\columnwidth - 4\tabcolsep) * \real{0.3000}}|
  >{\raggedright\arraybackslash}p{(\columnwidth - 4\tabcolsep) * \real{0.3000}}|
  >{\raggedright\arraybackslash}p{(\columnwidth - 4\tabcolsep) * \real{0.3000}}|}
\caption{Результаты тестовых проверок}\label{tab:test-results}\\
\hline
\begin{minipage}[b]{\linewidth}\raggedright
Проверка
\end{minipage} & \begin{minipage}[b]{\linewidth}\raggedright
Статус
\end{minipage} & \begin{minipage}[b]{\linewidth}\raggedright
Комментарий
\end{minipage} \\
\hline
\endhead
\hline
\endlastfoot
\passthrough{\lstinline!make test-unit!} & 112 тестов пройдено & Бизнес-логика шести сервисов \\ \hline
\passthrough{\lstinline!make test-e2e!} & 19 шагов пройдено & Сквозной пользовательский сценарий при поднятой платформе \\ \hline
\passthrough{\lstinline!make test-integration!} & 113 проверок пройдено & Все группы endpoint через API Gateway \\ \hline
\end{longtable}

\subsection{Матрица трассируемости требований}

После описания публичного API и уровней тестирования соответствие между функциональными требованиями, их реализацией и проверкой можно представить в явном виде. Каждое функциональное требование, сформулированное в аналитической главе, сопоставлено с реализующим его сервисом, ключевым endpoint публичного API и уровнем тестирования, на котором подтверждается его работоспособность (U --- unit, I --- интеграционный тест через API Gateway, E --- E2E-сценарий). Результат приведён в таблице~\ref{tab:traceability}.

\begin{longtable}[]{|
  >{\raggedright\arraybackslash}p{(\columnwidth - 8\tabcolsep) * \real{0.26}}|
  >{\raggedright\arraybackslash}p{(\columnwidth - 8\tabcolsep) * \real{0.20}}|
  >{\raggedright\arraybackslash}p{(\columnwidth - 8\tabcolsep) * \real{0.40}}|
  >{\centering\arraybackslash}p{(\columnwidth - 8\tabcolsep) * \real{0.06}}|}
\caption{Матрица трассируемости: требование --- сервис --- endpoint --- тест}\label{tab:traceability}\\
\hline
Требование & Сервис & Ключевой endpoint & Тест \\
\hline
\endfirsthead
\caption*{Продолжение таблицы \ref{tab:traceability}}\\
\hline
Требование & Сервис & Ключевой endpoint & Тест \\
\hline
\endhead
\hline
\endlastfoot
Регистрация и аутентификация, роли & Auth & \path|POST /api/v1/auth/register|, \path|/login|, \path|/refresh| & U, I, E \\ \hline
Разграничение прав по роли & Gateway, все сервисы & проверка JWT и роли в Gateway, заголовки \path|X-User-Role| & U, I \\ \hline
Заявки и связь тренер-спортсмен & User & \path|POST /api/v1/connections/request|, \path|.../accept| & U, I, E \\ \hline
Тренировочные группы & User & \path|POST /api/v1/groups|, \path|.../members| & U, I, E \\ \hline
Планы и назначения & Training & \path|POST /api/v1/training/plans| & U, I, E \\ \hline
Статусы и архивация назначений & Training & \path|GET /api/v1/training/assignments| & U, I \\ \hline
Шаблоны тренировок & Training & \path|/api/v1/training/templates| & U, I \\ \hline
Отчёты & Training & \path|POST /api/v1/training/assignments/{id}/report| & U, I, E \\ \hline
Уведомления & Notification & \path|GET /api/v1/notifications| & I, E \\ \hline
Аналитика & Analytics & \path|GET /api/v1/analytics/athletes/{id}/summary| & U, I \\ \hline
AI-рекомендации & AI & \path|POST /api/v1/ai/athletes/{id}/recommendations| & U, I \\ \hline
\end{longtable}

Матрица показывает, что все функциональные требования первой версии реализованы и покрыты тестами на уровне интеграционных проверок через API Gateway, а ключевые сценарии --- дополнительно unit- и E2E-тестами.

\subsection{Наблюдаемость}

Для того чтобы работоспособность платформы можно было не только проверить тестами, но и наблюдать количественно во время работы, все семь сервисов инструментированы метриками. Инструментирование реализовано единым middleware в разделяемом Go-модуле \passthrough{\lstinline!pkg/httpmetrics!}, который подключается к каждому сервису так же, как общий пакет событий. Каждый сервис экспортирует метрики в формате Prometheus на endpoint \passthrough{\lstinline!/metrics!}.

Собираются метрики модели RED (Rate, Errors, Duration): счётчик \passthrough{\lstinline!http\_requests\_total!} с метками сервиса, метода, маршрута и кода ответа, а также гистограмма задержек \passthrough{\lstinline!http\_request\_duration\_seconds!}. В качестве метки маршрута используется шаблон маршрута Echo (например, \path|/api/v1/training/assignments/:assignmentId|), а не сырой URL --- это ограничивает кардинальность временных рядов, не давая ей расти от идентификаторов в пути. В API Gateway middleware расположен до проверки JWT, поэтому запросы, отклонённые по авторизации, также попадают в метрики.

Сбор метрик выполняет Prometheus, который опрашивает endpoint \passthrough{\lstinline!/metrics!} всех сервисов с заданным интервалом и хранит полученные временные ряды. Для визуализации в конфигурацию развёртывания включена Grafana, которая поднимается вместе с платформой и автоматически подключает Prometheus как источник данных. Такое инструментирование не предназначено для нагрузочного тестирования прототипа, а закладывает основу для эксплуатационного мониторинга при дальнейшем развёртывании платформы: по метрикам видно распределение запросов по сервисам, доля ошибок и задержки обработки.

\subsection{Выводы по главе 3}

В третьей главе была рассмотрена реализация backend-платформы CoachLink. Реализованы сервисы, обеспечивающие регистрацию и авторизацию пользователей, управление профилями и связями, тренировочные группы, планы, назначения, отчёты, уведомления, аналитику и AI-рекомендации.

Система построена вокруг единого API Gateway, который проверяет JWT и маршрутизирует запросы. Межсервисное взаимодействие реализовано через HTTP и NATS JetStream. Данные разделены по сервисам, что соответствует выбранной микросервисной архитектуре.

Проведённая проверка показывает, что ключевая бизнес-логика покрыта unit-тестами, а сценарии уровня E2E и integration успешно выполняются при поднятой платформе. Это подтверждает работоспособность backend-платформы как на уровне отдельных сервисов, так и на уровне межсервисного взаимодействия через API Gateway.
